\documentclass[11pt,a4paper]{article}

\usepackage[T1]{fontenc}
\usepackage[utf8]{inputenc}
\usepackage[margin=2.4cm]{geometry}
\usepackage{amsmath,amssymb}
\usepackage{booktabs}
\usepackage{array}
\usepackage{xcolor}
\usepackage{enumitem}
\usepackage[colorlinks=true,linkcolor=blue!55!black,urlcolor=blue!55!black]{hyperref}
\usepackage{titlesec}

\definecolor{acc}{RGB}{40,90,160}
\definecolor{warn}{RGB}{190,120,20}
\titleformat{\section}{\large\bfseries\color{acc}}{\thesection.}{0.5em}{}
\titleformat{\subsection}{\normalsize\bfseries}{\thesubsection}{0.5em}{}
\setlist[itemize]{leftmargin=1.4em,itemsep=1pt,topsep=2pt}

\newcommand{\note}[1]{{\color{acc}\textbf{#1}}}

% simple boxed callout (standard packages only)
\newcommand{\warningbox}[2]{%
  \begin{center}
  \setlength{\fboxrule}{1pt}\setlength{\fboxsep}{9pt}%
  \fcolorbox{warn}{warn!7}{%
    \begin{minipage}{0.93\linewidth}
    \textbf{\color{warn}#1}\\[2pt]#2
    \end{minipage}}%
  \end{center}}

\title{\vspace{-1.2cm}\textbf{Weekly Progress Report}\\[2pt]
  \large FADO: Fairness-Aware Online Learning Under Concept Drift}
\author{}
\date{Week ending 3 July 2026}

\begin{document}
\maketitle
\vspace{-1.0cm}

\section*{Summary}
Good news this week: I extended \textsc{FADO} to \textbf{multiple protected attributes}
(intersectional fairness) --- both the code and the experiments to back it up. The
\textsc{Folktables} intersectional sweep is done, and the payoff is nice: the drift-reaction
controller now delivers a \emph{statistically significant} fairness gain over the
controller-free baseline --- something the old single-attribute setting was never able to show.
The \textsc{COMPAS} sweep isn't finished yet, for a slightly unusual reason (see the note below).

\warningbox{A small logistical hiccup --- experiments on pause}{%
Confession: my machines and the current heatwave are not on speaking terms. With no air
conditioning at home, running the PCs at full load for hours turns my room into a small sauna,
and the hardware isn't thrilled either. So sweeps and experiments are
\textbf{on hold} until either weather cools down (maybe). Everything already finished (the \textsc{Folktables} results below) is safe and sound
--- it's only the pending runs that are waiting for a less tropical week.}

\section{Multiple protected attributes (intersectional fairness)}

\subsection*{Motivation and design decision}
Previously \textsc{FADO} enforced fairness on a single binary attribute. We extended it to
enforce fairness jointly over \textbf{two} attributes. There were two candidate designs:
\begin{itemize}
  \item \textbf{Path A -- Intersectional.} Optimise demographic parity (DP) over the
        Cartesian-product subgroups: the joint group index is
        $a = 2\,a^{(1)} + a^{(2)} \in \{0,1,2,3\}$ ($K=4$ groups).
  \item \textbf{Path B -- Marginal sum.} Regularise each attribute independently.
\end{itemize}
We adopted \textbf{Path A}. Kearns et al.\ (ICML 2018, \emph{Preventing Fairness
Gerrymandering}) prove that marginal-only regularisation is \emph{provably defeatable}: a
model can achieve perfect per-attribute parity while being arbitrarily unfair on the joint
subgroups. Path A closes this loophole by construction, aligns with the dominant
intersectional-fairness definition (Foulds et al., ICDE 2020), and --- importantly --- is
also the smaller code change, since the node-level regulariser auto-scales to any $K$.

\subsection*{Implementation}
\begin{itemize}
  \item \textbf{Data loaders:} composite group index built in the loaders ---
        \textsc{COMPAS} uses \texttt{race}\,$\times$\,\texttt{sex},
        \textsc{Folktables} uses \texttt{SEX}\,$\times$\,\texttt{RAC1P} (binarised).
        Row-alignment invariants preserved through the per-seed subsampling and the
        \textsc{COMPAS} drift edits.
  \item \textbf{Metrics:} the intersectional DP (over the four joint cells) is the optimised
        objective; the two \emph{marginal} per-attribute DP values are computed and stored as
        diagnostics for every run.
  \item \textbf{Learning algorithm unchanged:} the fair oblique forest, the dual-ADWIN
        detector, and the reaction controller required \emph{no} modification --- they operate
        on whatever number of groups the data provides.
\end{itemize}

\section{Experiments}
All runs use the prequential (test-then-train, batch size 1) protocol, 30 random seeds,
paired Wilcoxon signed-rank tests with Holm--Bonferroni correction.
\begin{itemize}
  \item \note{Folktables (ACS Income, CA), sex\,$\times$\,race --- COMPLETE.}
        30 seeds $\times$ 4 models (FADO, controller-free Aranyani, ARF, RFR).
  \item \note{COMPAS, race\,$\times$\,sex --- PAUSED} (see availability note); the pipeline is
        implemented and validated on single seeds, only the full 30-seed sweep is outstanding.
\end{itemize}

\section{Results (Folktables, intersectional DP)}
Mean $\pm$ std over 30 seeds. DP is the intersectional violation over the four joint
subgroups (the objective \textsc{FADO} optimises); $\mathrm{DP}_{\text{sex}}$ and
$\mathrm{DP}_{\text{race}}$ are the marginal diagnostics. Lower DP is better; \textbf{bold}
marks the best mean per column.

\begin{center}
\renewcommand{\arraystretch}{1.2}
\begin{tabular}{l cccc}
\toprule
Model & Accuracy $\uparrow$ & DP $\downarrow$ & $\mathrm{DP}_{\text{sex}}\downarrow$ & $\mathrm{DP}_{\text{race}}\downarrow$ \\
\midrule
\textbf{FADO}            & $0.7979$          & $\mathbf{0.0688}$ & $\mathbf{0.0234}$ & $\mathbf{0.0243}$ \\
Aranyani (no controller) & $\mathbf{0.8005}$ & $0.0727$          & $0.0283$          & $0.0294$ \\
ARF                      & $0.7446$          & $0.1497$          & $0.0735$          & $0.0620$ \\
RFR                      & $0.7779$          & $0.1427$          & $0.0724$          & $0.0591$ \\
\bottomrule
\end{tabular}
\end{center}

\subsection*{Key findings}
\begin{itemize}
  \item \textbf{FADO more than halves the intersectional DP of both external baselines}
        ($\sim$54\% vs ARF, $\sim$52\% vs RFR). All twelve FADO-vs-\{ARF, RFR\} comparisons
        are significant at $p<0.001$ under \emph{both} the paired Wilcoxon and the Welch test.
  \item \textbf{The controller now beats the controller-free baseline.} FADO reduces
        intersectional DP by $0.38$pp over the bare Aranyani forest, significant under
        \emph{both} tests (Wilcoxon $p<0.001$, Welch $p<0.01$), with matching wins on both
        marginal axes. In the single-attribute setting this comparison was only a statistical
        tie --- so \textbf{the intersectional objective makes the controller's contribution
        measurable}, at a negligible accuracy cost ($-0.26$pp, not significant under Welch).
  \item \textbf{Gerrymandering confirmed empirically.} The ratio of intersectional to worst
        marginal DP is $\approx$2--3 for every method, and is \emph{highest} for the fair
        learners (FADO $\approx$2.8, Aranyani $\approx$2.5): they scrub the marginal axes so
        well that their residual unfairness is concentrated in the joint subgroups. A
        marginal-only audit would have declared FADO ``very fair'' while missing a
        joint-subgroup gap almost 3$\times$ larger --- direct empirical support for optimising
        the intersectional objective.
  \item \textbf{Statistics independently verified.} All significance values were cross-checked
        with two independent libraries (\texttt{scipy.stats} and
        \texttt{statsmodels.stats.multitest}); results match exactly.
\end{itemize}

\section{Code quality: controller audit}
I audited the full drift-reaction controller logic. Findings:
\begin{itemize}
  \item Removed a dead, expensive ($O(T^2)$) fairness computation from the inner loop
        (measurable speed-up, no change to results).
  \item Identified three discrepancies between the documented design and the executed code
        (recovery driven by cumulative rather than rolling accuracy; the warning/confirmation
        ADWIN thresholds are inverted; the label-noise guard is weaker than described).
  \item \textbf{Decision:} these were left \emph{unchanged} for now. On analysis, the current
        behaviour is plausibly what drives FADO's advantage on the abrupt-drift regime, so
        ``fixing'' them risks erasing the contribution; they are documented for a controlled
        follow-up rather than a blind change.
\end{itemize}

\section{Next steps (once the heatwave lets up)}
\begin{itemize}
  \item Finish the \textsc{COMPAS} intersectional 30-seed sweep.
  \item \textbf{Optimise equalised odds (EO) directly} (currently DP is the objective and EO is
        only monitored); the EO regulariser is implemented and driven by the same controller,
        so this is a configuration + sweep, then a second results block.
  \item Re-verify the gerrymandering and controller-gain claims hold on \textsc{COMPAS}.
\end{itemize}

\section*{A note on where I'd like to focus}
One thing I'd genuinely like to invest more time in is a proper \textbf{ablation study of the
reaction controller}. Right now it is the single weakest and most easily-targeted part of the
work: a reviewer can reasonably ask ``how much does each piece of the controller actually
contribute, and would a simpler variant do just as well?'' --- and at the moment I don't have a
clean, systematic answer. Isolating each component (the dual-ADWIN trigger, the warning /
confirmation / recovery phases, the rolling-window logic) and measuring its individual effect
would let me answer that head-on and turn a likely point of criticism into a strength. I'm
convinced this kind of ablation would meaningfully improve our odds at future conferences, even
though it is time-consuming and won't produce a headline result on its own. I'd appreciate the
room to spend a bit of extra time here rather than rushing straight to the next dataset.

\end{document}
